\documentclass[11pt]{article}

\usepackage[margin=1in]{geometry}
\usepackage[T1]{fontenc}
\usepackage[utf8]{inputenc}
\usepackage{hyperref}
\usepackage{amsmath}

\title{Outline: A Dual-Memory Patient World Model for Generative EHR}
\author{Project Working Draft}
\date{Last updated: 2026-04-12}

\begin{document}
\maketitle

\section{Status and Scope}

This file is a living manuscript scaffold for the repository's v2 line. It is
not a polished paper draft. Its purpose is to keep the implementation,
literature positioning, and scientific claim aligned while the codebase is
refactored toward a stronger conceptual contribution.

The active thesis is no longer ``hierarchical EHR transformer with better long
memory.'' The active thesis is a \textbf{dual-memory patient world model}:
\begin{itemize}
    \item typed fused events rather than raw token bundles
    \item marked time-to-event training rather than plain token prediction
    \item a compressive latent health state for current condition
    \item exact patient-internal memory for facts that should not be compressed away
    \item cohort-external precedent memory for analogous prior states and futures
    \item intervention-conditioned prospective rollout
\end{itemize}

\section{Working Title Options}

\begin{itemize}
    \item A Dual-Memory Patient World Model for Generative Electronic Health Records
    \item Generative EHR with Predictive-State Retrieval, Exact Patient Memory, and Marked Event Modeling
    \item A Typed Dual-Memory World Model for Longitudinal EHR Simulation
\end{itemize}

\section{One-Paragraph Thesis}

Electronic health records combine heterogeneous symbolic events, continuous
measurements, explicit care-regime transitions, and irregular timing over long,
multi-admission trajectories. Existing models typically flatten these records
into token or event streams and rely on one compressed context to summarize the
patient history. We instead propose a dual-memory patient world model for
structured EHR. Local attention operates over typed fused events inside semantic
care windows; a global latent health state represents current condition and
transition hazard; a patient-internal exact memory retains sparse personalized
facts that should not be compressed; and a cohort-external precedent memory
retrieves similar prior states together with their observed future
continuations. Training centers on marked time-to-event generation and
intervention-conditioned rollout. The intended contribution is not simply a new
backbone or tokenizer, but a reconceptualization of EHR generation as
\emph{state-based patient simulation} with both exact memory and analogical
future retrieval.

\section{Primary Claim Hierarchy}

\subsection{Primary claim}

\begin{quote}
Generative EHR modeling benefits from separating compressive state,
patient-specific exact memory, and cohort-level precedent memory, rather than
forcing all long-range information into either a single latent state or a long
token prefix.
\end{quote}

\subsection{Secondary claims}

\begin{itemize}
    \item EHR sequence modeling should compare patients by predictive clinical state, not only by token-prefix similarity.
    \item Numeric measurements should be modeled as continuous marks, not merely as extra tokens.
    \item Structural care transitions should be explicit modeling objects and natural window boundaries.
    \item Intervention-conditioned rollout is a more appropriate generative framing for clinical forecasting than unconstrained next-token continuation.
\end{itemize}

\subsection{Claims to avoid}

\begin{itemize}
    \item Do not claim that memory-augmented sequence modeling is new in general machine learning.
    \item Do not claim that retrieval by itself implies causal counterfactual validity.
    \item Do not claim that a richer recurrent block alone is the contribution.
    \item Do not claim a universal tokenizer for all EHR modalities.
\end{itemize}

\section{Central Reconceptualization}

The key reframing is:
\begin{quote}
EHR generation should operate on \emph{predictive clinical state}, not merely
on a tokenized prefix of observed history.
\end{quote}

The practical implication is that two patient histories may differ substantially
in observed token space while still being close in state space if they imply
similar distributions over plausible short- and medium-horizon futures under
similar intervention context.

This motivates three distinct state-bearing mechanisms:
\begin{itemize}
    \item a latent state for compressive current condition
    \item an internal exact memory for personalized facts
    \item an external precedent memory for analogical future support
\end{itemize}

This is the central scientific move. The target system is therefore not just a
forecaster, but a \emph{patient world model}. The closest formal intuition is a
predictive-state representation: boundary states should be compared by the
future behavior they imply, not only by how similar their past token prefixes
looked \cite{littman2001psr}.

\section{Novelty Positioning}

As of 2026-04-12, none of the following ingredients are novel in isolation:
\begin{itemize}
    \item continuous-time latent dynamics for irregular medical or event sequences
    \item attention-based or marked temporal point-process modeling
    \item explicit memory modules in general sequence modeling
    \item retrieval over large medical histories for discriminative prediction
\end{itemize}

The novelty should instead be framed as the combination:
\begin{itemize}
    \item typed fused EHR events with marked generative training
    \item semantic care-window hierarchy
    \item a compressive latent health state for current condition
    \item patient-internal exact memory for sparse personalized facts
    \item cohort-external precedent retrieval over predictive clinical state
    \item intervention-conditioned prospective rollout
\end{itemize}

The paper should explicitly state that memory-augmented generation is known in
general sequence modeling \cite{rae2020compressive,wu2022memorizing,bulatov2022rmt,borgeaud2022retro,behrouz2025titans},
and that retrieval over very long EHR histories is already useful in
discriminative settings \cite{kim2024remed}. The contribution is the
\emph{EHR-specific world-model decomposition} and the generative use of two
memory systems with a typed marked objective.

\section{Related Work Map}

\subsection{Hierarchical and long-context EHR models}

Hierarchical EHR transformers established that visit-level structure matters
when raw patient history exceeds the effective receptive field of a flat model
\cite{li2021hibehrt,odgaard2024corebehrt}. Linear-time EHR foundation models
such as EHRMamba strengthen the long-context baseline but still rely on a
single compressed state path rather than an explicit exact-memory mechanism
\cite{fallahpour2025ehrmamba}.

\subsection{Continuous-time latent modeling in health}

Latent ODEs established continuous-time latent dynamics for irregularly sampled
data \cite{rubanova2019latentode}. In health-specific settings, MIRA uses Neural
ODE extrapolation for medical time-series forecasting \cite{li2025mira}, while
TrajSurv learns continuous latent trajectories for survival modeling
\cite{zeng2025trajsurv}. These works motivate continuous or hybrid latent
dynamics, but they do not provide the dual-memory generative formulation.

\subsection{Marked event modeling}

Neural temporal point processes already provide a more faithful framework for
irregular event modeling than plain next-token prediction
\cite{enguehard2020ehrtpp,song2024dmtpp,liu2026nextpp}. ORA makes the
EHR-specific argument that pretraining should jointly model timing and marks
rather than rely on language-style token loss \cite{jing2026ora}. This project
inherits that objective stance.

\subsection{Memory-augmented sequence modeling}

General long-range sequence modeling has moved toward explicit memory:
compressed memory \cite{rae2020compressive}, non-differentiable exact memory
\cite{wu2022memorizing}, recurrent memory tokens \cite{bulatov2022rmt},
retrieval-augmented generation \cite{borgeaud2022retro}, and test-time memory
systems that separate long-term memory from attention
\cite{behrouz2025titans}. PromptTPP further shows that event-sequence models can
benefit from streaming retrieval without retraining \cite{xue2023prompttpp}.

\subsection{Retrieval in clinical and time-series modeling}

REMed demonstrates that near-infinite EHR history retrieval is powerful for
clinical prediction \cite{kim2024remed}. RAFT shows that retrieval of analogous
historical futures can improve forecasting in general time series
\cite{han2025raft}. These are especially relevant to the precedent-memory idea:
the value is not only remembering the current patient, but also retrieving what
happened to similar prior states.

\subsection{Clinical rollout and intervention modeling}

G-Transformer is a key antecedent for intervention-conditioned patient rollout
\cite{xiong2024gtransformer}. The present project should borrow the idea of
simulation under treatment regimes while staying conservative about causal
claims. Without stronger causal assumptions, the correct phrasing is
\emph{intervention-conditioned prospective rollout}, not unrestricted
counterfactual inference.

\subsection{Generative EHR baselines}

POPCORN is an important generative EHR point-process baseline
\cite{bhave2021popcorn}. It helps anchor the claim that the project is not
simply ``doing generative EHR for the first time,'' but rather proposing a
different generative factorization and memory architecture.

\section{Model Section Outline}

\subsection{Problem definition}

Represent the patient trajectory as an irregular sequence of typed events
\[
    \mathcal{E} = \{(t_i, f_i, c_i, u_i, m_i)\}_{i=1}^{N},
\]
where:
\begin{itemize}
    \item $t_i$ is event time
    \item $f_i$ is coarse family
    \item $c_i$ is concept identity
    \item $u_i$ is value or payload mark
    \item $m_i$ is a modifier set
\end{itemize}

Semantic windows partition the sequence into care regimes. Generation proceeds
over windows and events, not only over raw token prefixes.

\subsection{Typed event algebra and codecs}

\begin{itemize}
    \item Numeric measurements use the existing cVAE-plus-RVQ codec.
    \item Symbolic families use canonical concept identities plus modifiers.
    \item Structural transitions remain first-class events.
    \item All families map into a fused event interface before local modeling.
\end{itemize}

\subsection{Local event model}

Local attention runs over typed fused events within semantic windows. Chunking is
used only as a compute device. The local model has two responsibilities:
\begin{itemize}
    \item bind current within-window evidence for generation
    \item emit boundary summaries and memory-write candidates
\end{itemize}

\subsection{Global latent health state}

The latent is a belief state over current clinical condition:
\begin{align*}
    g_w^- &= \text{Drift}(g_{w-1}^+, \Delta \tau_w, \text{meta}_w) \\
    g_w^+ &= \text{Jump}(g_w^-, s_w, \text{type}_w),
\end{align*}
where $s_w$ is a compressive summary emitted by the local model.

The latent should encode:
\begin{itemize}
    \item current burden and instability
    \item support intensity
    \item unresolved problems carried forward
    \item transition hazard and tempo
\end{itemize}

The latent should not be treated as the patient file.

\subsection{Patient-internal exact memory}

The internal memory stores sparse exact facts that should not be compressed away.
Conceptually it is split into:
\begin{itemize}
    \item static memory
    \item persistent memory
    \item episodic memory
\end{itemize}

This memory is patient-specific, online, and exact. It is intended for:
\begin{itemize}
    \item durable chronic facts
    \item major procedures and support history
    \item critical rare events
    \item baseline non-competing facts
\end{itemize}

\subsection{Cohort-external precedent memory}

The second memory is the main conceptual step.

Each precedent item should contain:
\begin{itemize}
    \item a key representing a pre-window predictive clinical state
    \item care-regime and intervention metadata
    \item a compact future summary
    \item optionally an exact future continuation snippet
\end{itemize}

Retrieval should therefore answer:
\begin{quote}
Which prior patients were in a similar predictive clinical state, under similar
current intervention context, and what happened next?
\end{quote}

This differs sharply from token-prefix retrieval.

\subsection{Precedent item schema and future values}

The precedent index should store one item per semantic-window boundary, not one
item per token span. Each item should contain:
\begin{itemize}
    \item identity: subject, trajectory, boundary, and anchor-window ordinals
    \item anchor metadata: current window type, start time, duration, and previous-gap time
    \item coarse support/intervention flags
    \item a dense retrieval key built from the boundary packet, latent state, and a compact persistent-memory digest
    \item compact multi-horizon future summaries
    \item an optional exact continuation reference back into packed precompiled shards
\end{itemize}

The value side should be structured rather than token-level. A useful design is
three future horizons:
\begin{itemize}
    \item $h_1$: the next semantic window
    \item $h_2$: the next two windows or the next 24 hours
    \item $h_3$: the next four windows or the remainder of admission capped at 7 days
\end{itemize}

Each horizon should summarize the future with targets that are well supported
by MIMIC:
\begin{itemize}
    \item next-window type and timing
    \item event-family and payload histograms
    \item support or intervention flags
    \item transition flags
    \item event and measurement density
    \item extreme-measurement counts and numeric severity summaries
    \item medium-horizon disposition flags, with readmission targets as optional lower-priority additions
\end{itemize}

\subsection{Predictive-state similarity}

The paper should define state similarity functionally:

Two boundary states are similar if, under comparable current intervention and
regime context, they induce similar distributions over:
\begin{itemize}
    \item next window type
    \item next window timing
    \item short-horizon support evolution
    \item near-future event mixture
    \item medium-horizon disposition or readmission outcomes
\end{itemize}

This is the scientific justification for precedent retrieval.

\subsection{Generative loop}

The intended causal loop is:
\begin{enumerate}
    \item read internal patient memory using current latent and boundary metadata
    \item read external precedent memory using the current predictive state query
    \item run local attention conditioned on latent state, patient memory, precedent memory, and intervention context
    \item generate the current or next event/window
    \item update internal patient memory with exact writes
    \item update the latent state with a compressive summary
\end{enumerate}

\subsection{Marked generative objective}

The core objective remains:
\[
    \log p(\Delta t, f, c, u \mid \mathcal{H}, M_{\text{patient}}, M_{\text{precedent}}, a),
\]
where:
\begin{itemize}
    \item $\mathcal{H}$ is the observed history
    \item $M_{\text{patient}}$ is patient-internal exact memory
    \item $M_{\text{precedent}}$ is external precedent memory
    \item $a$ is current intervention context
\end{itemize}

Discrete marks use cross-entropy, continuous marks use family-conditioned
likelihood terms, and timing uses continuous-time likelihood.

\subsection{Intervention-conditioned rollout}

The paper should treat rollout as a first-class use case. The clean claim is:
\begin{itemize}
    \item the model supports prospective rollout under specified intervention context
    \item the model does not automatically justify unrestricted counterfactual claims
\end{itemize}

\section{Implementation Path}

The repo should be evolved toward the dual-memory world model in the following
order:
\begin{enumerate}
    \item introduce a canonical boundary-state packet as the interface between local attention and all long-range modules
    \item refactor patient-internal memory into explicit static, persistent, and episodic banks
    \item build a cohort-external precedent index over boundary states with aligned multi-horizon future summaries and exact continuation references
    \item add a future-aware retrieval objective so precedent similarity reflects future agreement rather than token-prefix overlap
    \item condition generation on latent state, patient memory, and precedent memory jointly
    \item only then finalize the latent-state family
\end{enumerate}

The first engineering principle is therefore:
\begin{quote}
Do not finalize the latent before the boundary-state packet and dual-memory
interfaces are fixed.
\end{quote}

This is important because the current single-vector summary and the current
single-vector latent are migration bridges, not the final conceptual objects.

\section{Why This Is More Than ``A Better Latent''}

The paper should be explicit here. A stronger recurrent or continuous-time
latent backbone alone is not the main contribution. The main contribution is the
division of labor:
\begin{itemize}
    \item latent state for current condition
    \item exact internal memory for patient-specific facts
    \item precedent memory for analogical future support
\end{itemize}

Once that contract is fixed, the final latent family can be chosen more
cleanly. Before that, a more expressive latent mostly risks becoming an
accidental memory store.

\section{Experiments Section Outline}

\subsection{Baselines}

\begin{itemize}
    \item v1 project baseline
    \item hierarchical transformer baselines such as Hi-BEHRT and CORE-BEHRT
    \item long-context/state-space baselines such as EHRMamba
    \item generative event baselines such as POPCORN
    \item retrieval baselines such as REMed-style or RAFT-style retrieval without dual-memory decomposition
    \item intervention-rollout comparison to G-Transformer-style settings where applicable
\end{itemize}

\subsection{Core ablations}

\begin{itemize}
    \item token-path vs fused-event-path
    \item plain autoregression vs marked objective
    \item no memory vs patient memory only vs precedent memory only vs dual memory
    \item no latent vs latent only vs latent plus patient memory vs latent plus dual memory
    \item generic similarity vs future-aware predictive-state similarity
    \item latent family comparison once the dual-memory interface is fixed
\end{itemize}

\subsection{Task families}

\begin{itemize}
    \item next-event and next-window generative forecasting
    \item support escalation and deterioration prediction
    \item transfer, discharge, and readmission timing
    \item intervention-conditioned prospective rollout
    \item rare-pathway generation under long historical context
\end{itemize}

\subsection{Key diagnostics}

\begin{itemize}
    \item long-horizon degradation with and without each memory system
    \item retrieval quality measured by future agreement, not only embedding similarity
    \item behavior on patients with divergent token histories but similar future trajectories
    \item performance stratified by inter-admission gap and trajectory length
    \item reliance on patient memory vs precedent memory across clinical regimes
\end{itemize}

\section{Tables and Figures To Reserve Early}

\subsection{Figures}

\begin{itemize}
    \item overall dual-memory world-model architecture
    \item typed event algebra and payload-codec schematic
    \item patient memory vs precedent memory vs latent state
    \item predictive-state retrieval diagram
    \item intervention-conditioned rollout diagram
\end{itemize}

\subsection{Tables}

\begin{itemize}
    \item literature positioning table
    \item modality routing table
    \item ablation table
    \item long-horizon robustness table
    \item memory write/read policy table
    \item precedent retrieval evaluation table
\end{itemize}

\section{Literature Positioning Table Skeleton}

\begin{center}
\begin{tabular}{p{0.18\linewidth}p{0.13\linewidth}p{0.14\linewidth}p{0.14\linewidth}p{0.14\linewidth}p{0.18\linewidth}}
\hline
Model / paper & Hierarchical local/global & Continuous-time latent & Internal exact memory & External precedent retrieval & Generative / rollout framing \\
\hline
Hi-BEHRT \cite{li2021hibehrt} & yes & no & no & no & no \\
CORE-BEHRT \cite{odgaard2024corebehrt} & no & no & no & no & no \\
EHRMamba \cite{fallahpour2025ehrmamba} & no & no & no & no & partly \\
POPCORN \cite{bhave2021popcorn} & no & yes-like TPP timing & no & no & yes \\
REMed \cite{kim2024remed} & no & no & partial history retrieval & no & no \\
G-Transformer \cite{xiong2024gtransformer} & no & rollout-focused & no & no & yes \\
TrajSurv \cite{zeng2025trajsurv} & no & yes & no & no & no \\
ORA \cite{jing2026ora} & objective only & no & no & no & partly \\
Ours & yes & yes or hybrid & yes & yes & yes \\
\hline
\end{tabular}
\end{center}

\section{Open Risks}

\begin{itemize}
    \item External precedent retrieval may help only if the similarity space is future-aware rather than syntax-aware.
    \item The internal and external memories may collapse onto similar roles unless their write/read contracts are explicit.
    \item Intervention-conditioned rollout is scientifically strong but must not be oversold as causal counterfactual prediction.
    \item The final latent design should be chosen only after the dual-memory interface is fixed, otherwise architecture search may optimize the wrong object.
\end{itemize}

\bibliographystyle{plain}
\bibliography{foundation_model_v2_refs}

\end{document}
